\pagestyle{plain}
\chapter*{Resumen}

En este trabajo se describe y se presenta una solución de mejora de eficiencia temporal en el acceso a 
nubes de puntos 3D utilizando como estructura de datos \textit{octrees} lineales. El estudio consiste en el reordenamiento espacial
 de puntos para tratar de explotar la capacidad de poda a la hora de encontrar vecindades. Se evalúan resultados temporales usando diferentes técnicas de reordenamiento, operando tanto con las coordenadas cartesianas como con las coordenadas polares de los puntos. Los resultados experimentales muestran un impacto positivo en la eficiencia temporal bajo determinadas circunstancias, a cambio de un gasto en memoria utilizada.


